\section{The fundamental conjecture: three different comparison problems}\label{sec:E}

\subsection{The change of definition in \S29}

We use the conventions of Part A.  In particular, $\pi_{g,n}$ is the
\emph{profinite} surface group, $P_{g,n}$ is the pure orientation-preserving
mapping class group, and $\Sigma_{g,n}$ is the closed image of the full
orientation-preserving mapping class group, with permutations of the marked
points allowed, in
$O_{g,n}=\operatorname{Out}_{\mathrm{lac}}(\pi_{g,n})$.
The suffix ``lac'' retains the conjugacy classes of peripheral procyclic
subgroups and their common cyclotomic exponent.  All normalizers, images,
quotients and commutator closures in this section are topological.  An
identification with an automorphism group is an identification of continuous
automorphisms.  Choices of a geometric marking identify fundamental groups;
changing paths conjugates the resulting actions.  The distinction between a
canonical outer action and a chosen genuine action will be retained.

\begin{center}
\begin{tabular}{L{.20\linewidth}L{.45\linewidth}L{.26\linewidth}}
Our notation & Definition & Manuscript source\\ \hline
$N^{\mathrm{norm}}_{g,n}$ &
$N_{O_{g,n}}(\Sigma_{g,n})$ & \S26, fol.~164, scan~324, (14)\\
$\Gamma^{\mathrm{norm}}_{g,n}$ &
$N^{\mathrm{norm}}_{g,n}/\Sigma_{g,n}$ & same passage; fol.~168, scan~332, (27)\\
$M^{\mathrm{arith}}_{g,n}$ & closed image of the arithmetic universal-curve
fundamental group in $\operatorname{Aut}_{\mathrm{lac}}(\pi_{g,n})$ &
\S29, fols.~187--188, scans~370--371, (9)--(13)\\
$N^{\mathrm{arith}}_{g,n}$ & exterior image of $M^{\mathrm{arith}}_{g,n}$ & same passage\\
$\Gamma^{\mathrm{arith}}_{g,n}$ &
$N^{\mathrm{arith}}_{g,n}/\Sigma_{g,n}$ & same passage\\
$G_{\mathbb Q}$ & $\operatorname{Gal}(\overline{\mathbb Q}/\mathbb Q)$ &
his $\mathbb I\Gamma_{\mathbb Q}$\\
\end{tabular}
\end{center}

The symbols $N$ and $\mathbb I\Gamma$ do not keep their definitions across
these passages.  Scan~370 announces the change to an arithmetic construction;
on scan~371 we read the instruction to denote the \emph{image} of the
arithmetic homomorphism by $M$, and the arrow
from $G_{\mathbb Q}$ to the resulting quotient is explicitly called
surjective.  The latter conjecture asks for injectivity.  We checked these
two scans, including the words surrounding the formulas.  The missing
orientation superscript in one occurrence on scan~371 is recorded in the
readings register; it does not license identifying the two normalizers.

\statement{E.1}{Proved}{The early and revised conjectures are inequivalent statements}
For every hyperbolic type $2g-2+n>0$ there are continuous homomorphisms
\[
 G_{\mathbb Q}\twoheadrightarrow\Gamma^{\mathrm{arith}}_{g,n}
 \hookrightarrow\Gamma^{\mathrm{norm}}_{g,n}.
\]
The first is surjective by definition.  Its injectivity is the revised
conjecture of \S29, fol.~188, scan~371.  The conjecture of \S26,
fol.~168, scan~332, asks that the composite be an isomorphism; it additionally
asserts that every element of the normalizer quotient is arithmetic.
The stronger assertion $N^{\mathrm{norm}}_{g,n}=O_{g,n}$ appearing in that
passage is a further assertion, not an equivalent formulation of either one.
\begin{proof}
The geometric subgroup is normal in the arithmetic fundamental group.
Its image $\Sigma_{g,n}$ is consequently normalized by the arithmetic image.
Thus $N^{\mathrm{arith}}_{g,n}\subseteq N^{\mathrm{norm}}_{g,n}$, and both
quotients are by the same closed subgroup.  The exact sequence of the
arithmetic moduli stack gives the first surjection.  Surjectivity of the
composite is exactly equality of the two images modulo $\Sigma_{g,n}$,
equivalently equality of the two groups $N$.  No assertion about the rest of
$O_{g,n}$ follows.  This also explains why the discussion of admissible
splittings on \S33, fol.~242, scan~479, cannot be read using the original
normalizer definition without an additional argument.
\end{proof}

\statement{E.2}{Known}{The revised conjecture is true for every hyperbolic type}
The natural map
\[
 G_{\mathbb Q}\xrightarrow{\ \sim\ }\Gamma^{\mathrm{arith}}_{g,n}
\]
is an isomorphism of profinite groups for every hyperbolic $(g,n)$.
It is compatible with the arithmetic maps between types wherever those
maps are defined.  This includes $(0,3)$, $(0,4)$, $(0,5)$, $(1,1)$ and
$(1,2)$, as well as proper curves of genus at least two.  It does not require
the congruence subgroup property in unknown genera.
\begin{proof}
Part D, Statement D.17, gives the precise application of
\citet[Corollary~6.4]{HoshiMochizuki2011}: the arithmetic Galois image
has trivial intersection with the geometric universal-monodromy image.
After a finite field extension admitting a curve of the specified type,
this proves injectivity on an open subgroup of $G_{\mathbb Q}$.  The remaining
kernel is a finite normal subgroup of the slim group $G_{\mathbb Q}$, and
is therefore trivial.  Statement E.1 supplies surjectivity.  The use of the
actual image $\Sigma$, rather than an unproved identification of that image
with the full profinite completion, is essential.  Naturality follows from
the arithmetic fundamental-group diagrams.  A continuous bijection between
profinite groups is a homeomorphism.
\end{proof}

For $(0,3)$, the arithmetic moduli stack with unlabelled points is
$B\mathfrak S_3$ over $\mathbb Q$; the labelled stack is $\operatorname{Spec}
\mathbb Q$.  Belyi faithfulness supplies the required injectivity here:
the three-branch-point theorem is \citet[\S4, Theorem~4]{Belyi1980},
and faithfulness is its unnumbered corollary on printed page~256.
Consequently the impression recorded on scan~479 is correct for the
\emph{revised} group.  Lead~1 is false as an identification of that group
with an independently defined $\widehat{\mathrm{GT}}$: surjectivity onto
the revised arithmetic quotient is already built into its definition.

\statement{E.3}{Known}{The abelianization of the revised group}
For every hyperbolic type,
\[
 (\Gamma^{\mathrm{arith}}_{g,n})^{\mathrm{ab}}
 =\Gamma^{\mathrm{arith}}_{g,n}/
   \overline{[\Gamma^{\mathrm{arith}}_{g,n},\Gamma^{\mathrm{arith}}_{g,n}]}
 \xrightarrow{\ \sim\ }\widehat{\mathbb Z}^{\times}
\]
is the cyclotomic character.  Thus the consequence anticipated in
\S26--\S29 (in particular scan~340) is a theorem for the revised groups.
\begin{proof}
Apply E.2 and the Kronecker--Weber theorem:
$\mathbb Q^{\mathrm{ab}}=\bigcup_m\mathbb Q(\mu_m)$.
The restriction map identifies the indicated quotient with
$\varprojlim_m(\mathbb Z/m\mathbb Z)^{\times}$.  The image of inertia
generators in the universal curve is acted on by this same character.
This argument says nothing about additional abelian quotients of a larger
normalizer group.
\end{proof}

\subsection{The genus-zero comparison, including the isolated tripodal case}

We write $\widehat F_2=\langle x,y\rangle$, $z=(xy)^{-1}$.  A profinite
Grothendieck--Teichm\"uller element is represented by $(\lambda,f)$ with
$\lambda\in\widehat{\mathbb Z}^{\times}$ and
$f\in\overline{[\widehat F_2,\widehat F_2]}$, such that
\[
 x\longmapsto x^\lambda,\qquad
 y\longmapsto f^{-1}y^\lambda f
\]
defines an automorphism and satisfies Drinfeld's symmetry, hexagon and
pentagon equations.  The automorphism requirement is part of the definition
in the profinite setting.  We denote by $\widehat{\mathrm{GT}}_0$ the group
with symmetry and hexagon only.  The pentagon is a condition in the
fundamental group of $\mathcal M_{0,5}$, not an extra relation in
$\pi_1(\mathcal M_{0,4})$ with no reference to its embeddings.  Our convention
is that $P_{0,m}=\pi_1(\mathcal M_{0,m}(\mathbb C))$.

\statement{E.4}{Known}{The exact two-level theorem}
Let $\operatorname{Out}^{\#}(\widehat P_{0,m})$ consist of outer
automorphisms which send each standard boundary generator $x_{ij}$ to a
conjugate of $x_{ij}^{\lambda}$ for a single
$\lambda\in\widehat{\mathbb Z}^{\times}$, and which commute with the
outer relabelling action of $\mathfrak S_m$.  Then
\[
 \operatorname{Out}^{\#}(\widehat P_{0,4})
       \cong\widehat{\mathrm{GT}}_0,
 \qquad
 \operatorname{Out}^{\#}(\widehat P_{0,m})
       \cong\widehat{\mathrm{GT}}\quad(m\geq5).
\]
The forgetful maps identify the latter groups for all $m\geq5$, and the
map to the four-point group is the usual injection
$\widehat{\mathrm{GT}}\hookrightarrow\widehat{\mathrm{GT}}_0$.
\begin{proof}
These are precisely the Main Theorem (a),(b) and its Corollary in
\citet{HarbaterSchneps2000}; the four-point identification is their
Theorem~2, and the five-point identification their Theorem~3.
Their Proposition~7 proves the pentagon extension criterion, and
Proposition~8 proves injectivity of forgetting.  The common exponent and
the relabelling centralizer occur in their definition of
$\operatorname{Out}^{\#}$ and cannot be dropped from the statement.
\end{proof}

There is also a version with chosen genuine braid-group automorphisms:
the objects are $\widehat B_j$ and the subgroups generated by
$\sigma_1^2,\sigma_2,\ldots,\sigma_{j-1}$, with the standard inclusions
and doubling maps.  Automorphisms act trivially on the permutation
quotients and commute with those maps.  The tower through three strands
has automorphism group $\widehat{\mathrm{GT}}_0$; through four strands,
and at every further finite truncation, it has
$\widehat{\mathrm{GT}}$.  These are the tower statements recalled and
used in \citet[\S0.3]{HarbaterSchneps2000}.  They specify a concrete
meaning of ``two levels''.  Arbitrary independent outer automorphisms of
two groups are not such a tower action.

\statement{E.5}{Known}{The intrinsic five-point and higher automorphism theorem}
For $m\geq5$,
\[
 \operatorname{Out}(\widehat P_{0,m})
       \cong\widehat{\mathrm{GT}}\times\mathfrak S_m.
\]
The symmetric factor is the geometric relabelling action; the other factor
preserves all generalized fibre subgroups.  Both factors and their
commutation are part of the theorem.
\begin{proof}
Take $(g,r)=(0,3)$ and $n=m-3\geq2$ in
\citet[Corollaries~2.6(i),(iii) and~2.8(i),(ii)]{HoshiMinamideMochizuki2022}.
Their configuration group $\Pi_n$ is $\widehat P_{0,m}$; their
$\operatorname{Out}^{gF}$ is the factor displayed here.  Their
Definition~2.7 identifies $\mathrm{GT}^{\Sigma}$ with the classical
$\widehat{\mathrm{GT}}$ when the set of primes $\Sigma$ is all primes.
In particular the result does not presuppose that a given automorphism
preserves inertia.  It reconstructs the relevant distinguished subgroups.
\end{proof}

\statement{E.6}{Proved}{A normalizer lemma for the universal curve}
Let $\pi$ be a centre-free profinite group, let
$S^!\triangleleft S\subseteq\operatorname{Out}(\pi)$ be closed subgroups,
and let $E$ be the inverse image of $S^!$ in $\operatorname{Aut}(\pi)$.
If $S\subseteq N\subseteq\operatorname{Out}(\pi)$ and $N$ normalizes
$S$ and $S^!$, conjugation induces an injection
\[
 N/S^!\hookrightarrow\operatorname{Out}(E)
\]
whose image normalizes the image of $S/S^!$.
\begin{proof}
Choose a representative $a\in\operatorname{Aut}(\pi)$ of an element of
$N$.  Conjugation by $a$ preserves $E$.  Replacing $a$ by an inner
automorphism of $\pi$ changes this automorphism of $E$ by an inner
automorphism of $E$, since $\operatorname{Inn}(\pi)\subseteq E$.
Suppose conjugation by $a$ on $E$ equals conjugation by $e\in E$.
Then $e^{-1}a$ centralizes $\operatorname{Inn}(\pi)$.  For every
$u\in\pi$ we have
$\operatorname{inn}((e^{-1}a)(u))=\operatorname{inn}(u)$; centre-freeness
gives $(e^{-1}a)(u)=u$.  Hence $a=e$, and its exterior class belongs to
$S^!$.  This proves injectivity.  The normalizer assertion follows by
conjugating the subgroup $S$ before passage to the quotient.
\end{proof}

\statement{E.7}{Proved}{The early normalizer quotient in genus zero}
With the full permutation convention of \S19 and \S26,
\[
 \Gamma^{\mathrm{norm}}_{0,3}\cong\widehat{\mathrm{GT}}_0,
 \qquad
 \Gamma^{\mathrm{norm}}_{0,n}\cong\widehat{\mathrm{GT}}
       \quad(n\geq4).
\]
The identifications respect the arithmetic embeddings after the usual
geometric markings.  The arithmetic subgroup in every case is
$\Gamma^{\mathrm{arith}}_{0,n}\cong G_{\mathbb Q}$.
\begin{proof}
First suppose $n\geq4$.  The profinite Birman sequence, and genus-zero
congruence subgroup property from Parts B and C, identify the inverse
image of $\Sigma^!_{0,n}$ in $\operatorname{Aut}(\pi_{0,n})$ with
$E=\widehat P_{0,n+1}$.  A normalizer element preserves $\Sigma^!$ because
it preserves the peripheral permutation map.  E.6 embeds
$N^{\mathrm{norm}}/\Sigma^!$ into
$\operatorname{Out}(E)=\widehat{\mathrm{GT}}\times\mathfrak S_{n+1}$.
The image of $\Sigma/\Sigma^!$ is the point stabilizer
$\mathfrak S_n\subset\mathfrak S_{n+1}$.  Its normalizer is itself:
its unique fixed letter must be preserved by any normalizing permutation.
Consequently $N^{\mathrm{norm}}/\Sigma$ embeds in
$\widehat{\mathrm{GT}}$.

Conversely, the standard genus-zero tower action of
$\widehat{\mathrm{GT}}$ preserves the kernel of
$\widehat P_{0,n+1}\to\widehat P_{0,n}$ and commutes, outerly, with
the $\mathfrak S_n$ action.  That kernel is $\pi_{0,n}$.  Restriction
therefore produces an exterior automorphism of $\pi_{0,n}$ preserving
its peripheral classes and normalizing $\Sigma_{0,n}$.  An inner change
of a representative of the automorphism of $E$ changes this restriction
by $\Sigma^!$, and hence has no effect modulo $\Sigma$.
Applying E.6 shows that this construction is inverse to the injection.

For $n=3$, $\Sigma_{0,3}\cong\mathfrak S_3$ and
$\Sigma^!_{0,3}=1$.  Every automorphism of $\mathfrak S_3$ is inner.
Thus its normalizer in $O_{0,3}$ is the product of $\mathfrak S_3$ with
its centralizer, with trivial intersection.  The centralizer acts
trivially on the three peripheral labels.  Under
$\pi_{0,3}\cong\widehat P_{0,4}$, it is exactly
$\operatorname{Out}^{\#}(\widehat P_{0,4})$: the relabelling action of
$\mathfrak S_4$ on four-point moduli has image $\mathfrak S_3$.
Now apply E.4.  Naturality of the arithmetic monodromy gives the last
assertions.
\end{proof}

Replacing the full group $\Sigma_{0,3}$ by its pure subgroup $1$ changes
the problem: the pure normalizer quotient is then the entire
label-preserving inertia automorphism group of $\widehat F_2$.
It cannot be called $\widehat{\mathrm{GT}}$ merely because the same free
group occurs.  This distinction is especially relevant to retrospective
accounts which suppress the exclamation-mark convention of \S19.

\statement{E.8}{Open}{The remaining profinite arithmetic equality}
The assertion
\[
 G_{\mathbb Q}\longrightarrow\widehat{\mathrm{GT}}
       \quad\hbox{is surjective}
\]
is the extra arithmetic assertion needed to identify the early normalizer
quotients in types $(0,4)$ and $(0,5)$ with $G_{\mathbb Q}$.
The map is a continuous closed embedding by Belyi faithfulness and the
Galois compatibility of the tower.  The isolated $(0,3)$ assertion instead
asks for surjectivity onto $\widehat{\mathrm{GT}}_0$; it contains, in
addition, the requirement that every such automorphism satisfy the
pentagon.  We do not infer either surjectivity from E.2.

Our source audit did not verify a primary proof of a \emph{strict}
profinite inclusion $\widehat{\mathrm{GT}}\subsetneq
\widehat{\mathrm{GT}}_0$.  Some historical discussions assert such
strictness; we have not substituted a rational or a finite-shadow
calculation for the missing profinite counterexample.  Accordingly we
record the isolated-level equality with its precise additional condition,
rather than declare it false on that basis.  This is a limitation of this
audit, not a claim that no relevant theorem exists.  The invalid step in
Lead~1 is already settled independently by E.1--E.2.

\subsection{The two elliptic types and the meaning of an all-genus tower}

\statement{E.9}{Known}{The two-point elliptic automorphism theorem}
There is an isomorphism
\[
 P_{1,2}\cong B_4/Z(B_4),\qquad
 \operatorname{Out}(\widehat P_{1,2})\cong\widehat{\mathrm{GT}}.
\]
Here $P_{1,2}$ fixes both marked points individually.  The assertion is
\citet[Theorem~B, Theorem~4.3 and Corollary~C]{MinamideNakamura2022}.
It does not identify the outer automorphism group of
$\widehat{\mathrm{SL}_2(\mathbb Z)}$ with
$\widehat{\mathrm{GT}}$, nor the unrestricted group with permuted
elliptic markings with its pure counterpart.  For the full group in type
$(1,2)$ there is a central involution exchanging the two punctures.
The inertia-preserving outer group and the unrestricted outer group are
respectively $\widehat{\mathrm{GT}}$ and
$C_2\times\widehat{\mathrm{GT}}$ in
\citet[Proposition~3.9]{Boggi2026Automorphisms}.

\statement{E.10}{Proved}{The one-point elliptic normalizer}
There is an isomorphism
\[
 \Gamma^{\mathrm{norm}}_{1,1}\cong\widehat{\mathrm{GT}},
\]
compatible with the arithmetic inclusion.  This statement concerns the
normalizer in $\operatorname{Out}_{\mathrm{lac}}(\pi_{1,1})$, not all
outer automorphisms of the elliptic mapping class group.
\begin{proof}
The profinite Birman sequence and genus-one congruence theorem identify
the inverse image of $\Sigma_{1,1}$ in
$\operatorname{Aut}(\pi_{1,1})$ with $E=\widehat P_{1,2}$.
E.6 and E.9 give an injection of the asserted normalizer quotient into
$\widehat{\mathrm{GT}}$.

We verify the converse, including the peripheral condition.  Put
$\zeta_j=(\sigma_1\cdots\sigma_{j-1})^j$.  The homomorphism
$\beta:B_4\to B_3$, given by
$\sigma_1,\sigma_2,\sigma_3\mapsto\sigma_1,\sigma_2,\sigma_1$,
sends $\zeta_4$ to $\zeta_3^2$.  It induces the elliptic forgetful map
\[
 E=\widehat{B_4/\langle\zeta_4\rangle}
   \longrightarrow
 \widehat{B_3/\langle\zeta_3^2\rangle}
   =\widehat{\mathrm{SL}_2(\mathbb Z)},
\]
with kernel $\pi_{1,1}$.  Compose with the quotient to
$\widehat{B_3/\langle\zeta_3\rangle}$.  Its kernel $K$ is
$C_2\mathbin{\widehat *}C_2\mathbin{\widehat *}C_2$, and it is
characteristic in $E$; these are the Cardano--Ferrari calculation and
Proposition~3.2(iii) of \citet{MinamideNakamura2022}.
The subgroup $\pi_{1,1}$ is the kernel of the unique homomorphism
$K\to C_2$ which is nonzero on each conjugacy class of involutions.
Every finite-order element of this profinite free product is conjugate
into a factor.  Thus this kernel is characteristic in $K$, hence in $E$.

Under the Birman identification its peripheral element is, up to inverse
and conjugacy, $(\sigma_1\sigma_2)^6=\zeta_3^2$ in $E$: this is the
two-chain relation, in which the boundary of a one-holed torus also
bounds the disc containing the two marked points.  The standard
Grothendieck--Teichm\"uller braid action sends $\zeta_3$ to
$\zeta_3^\lambda$ and $\zeta_4$ to $\zeta_4^\lambda$; consequently
it preserves this peripheral subgroup with exponent $\lambda$.
Restriction to the characteristic subgroup $\pi_{1,1}$ therefore
defines an element of $O_{1,1}$ normalizing $\Sigma_{1,1}$.
The ambiguity from an inner automorphism of $E$ is exactly $\Sigma_{1,1}$
after exterior restriction.  E.6 shows that restriction and conjugation
are inverse.  This proves the claim.  The full-twist formulas used here
are Drinfeld's standard braid action, also stated explicitly in
\citet[Theorem~2.1]{LochakNakamuraSchneps2026}.
\end{proof}

\statement{E.11}{Proved}{An exact reduction for the two-point normalizer}
Put $E=\widehat P_{1,3}$, and let
$\pi=\ker(E\to\widehat P_{1,2})$ be the kernel of forgetting the third
point.  Let $A\subseteq\operatorname{Out}(E)$ consist of the classes
which preserve $\pi$, whose restrictions preserve its two peripheral
conjugacy classes with a common exponent, and which normalize the
outer $\mathfrak S_2$ action permuting the first two marks.  Then
\[
 \Gamma^{\mathrm{norm}}_{1,2}\cong A/\mathfrak S_2.
\]
In particular
$G_{\mathbb Q}\cong\Gamma^{\mathrm{arith}}_{1,2}
\hookrightarrow A/\mathfrak S_2$.
\begin{proof}
Use E.6 with $S^!=\Sigma^!_{1,2}$ and $S=\Sigma_{1,2}$.
The pullback of $S^!$ in $\operatorname{Aut}(\pi)$ is $E$ by the
Birman sequence and the genus-one congruence theorem.  Conversely, a
representative $b$ of a class in $A$ restricts to an automorphism
$a$ of $\pi$.  For $e\in E$ the identity
$b(eue^{-1})=b(e)a(u)b(e)^{-1}$, $u\in\pi$, implies that
$b$ on $E\subseteq\operatorname{Aut}(\pi)$ is conjugation by $a$.
The normalizer and peripheral conditions are therefore exactly the
conditions that its outer class belong to $N^{\mathrm{norm}}_{1,2}$.
Changing $b$ by an inner automorphism changes $a$ by $\Sigma^!$.
Quotienting by the remaining permutations proves the formula.
\end{proof}

This is an exact reduction, not an identification of the unrestricted
group $A/\mathfrak S_2$ with $\widehat{\mathrm{GT}}$.
The theorem E.9 concerns the base group $P_{1,2}$; the normalizer problem
also retains its universal curve, whose group is $P_{1,3}$.
Applying E.9 to the wrong group would erase precisely this distinction.

\statement{E.12}{Known}{Refined actions with quilt and clutching data}
Let $\mathcal G_{\mathrm{NS}}$ denote the subgroup called
$\mathrm I\!\Gamma$ in \citet[Definition~1.1]{NakamuraSchneps2000},
defined by their additional relations $(\mathrm{III}')$ and $(\mathrm{IV})$,
with their Kummer cocycle $\rho_2$.  Then
\[
 G_{\mathbb Q}\hookrightarrow\mathcal G_{\mathrm{NS}}
       \hookrightarrow\widehat{\mathrm{GT}}.
\]
The first inclusion and subgroup assertion are their Theorem~1.2.
Their Theorems~1.3--1.4 give genuine actions after choosing a quilted
pants decomposition, canonical outer actions after forgetting that
choice, and compatibility with inclusions of subsurfaces cut out by
pants.  The change-of-quilt ambiguity is an explicitly described inner
automorphism.  The group $\Lambda$ of
\citet{HatcherLochakSchneps2000} is a related subgroup of the
$\lambda=1$ group, defined using the genus-one relation $(R)$; it is not
the full arithmetic Galois group.  Their Theorem~A (Lemma~7 of \S3)
proves the subgroup assertion; Theorem~B (Theorem~4 of \S3) gives the
pants-dependent genuine action and choice-independent outer action;
Theorem~C (Theorem~5 of \S3) embeds it in the compatible outer
automorphism group of their full subsurface tower.  These statements require the specified
extra compatibility data.  No equality between these named refined
subgroups and $\widehat{\mathrm{GT}}$ follows from E.9 or E.10.

\statement{E.13}{Known}{The current outer-tower results and their precise scope}
We distinguish two functors on orientable hyperbolic surfaces with
boundary and embeddings which may fill punctures:
\[
 S\longmapsto\widehat{P(S,\partial S)},\qquad
 S\longmapsto\check P(S,\partial S).
\]
Their arrows are the induced continuous homomorphisms modulo inner
automorphisms of the target.  For these specific functors,
\citet[\S2.8, Theorem~6.23 and Corollary~6.24]{Boggi2026Automorphisms}
proves
\[
 \operatorname{Aut}(\check{\mathfrak T}^{\mathrm{out}})
       \cong\widehat{\mathrm{GT}},\qquad
 \widehat{\mathrm{GT}}\hookrightarrow
       \operatorname{Aut}(\widehat{\mathfrak T}^{\mathrm{out}}).
\]
We use the explicitly dated version~6, 22 April 2026, as a preprint
theorem; we did not audit its proof in full.  Equality in the second
display remains its Conjecture~6.21.  The full profinite congruence
subgroup property would imply that equality.

These are statements about natural automorphisms of a group-valued
outer functor.  They are not, without a comparison theorem, statements
about all base-point paths, quilt moves, finite-cover maps or all
degeneration bimodules in a finer tower.  In particular they do not
silently remove the named relations in E.12.  The independent
\citet[Theorem~1.1]{LochakNakamuraSchneps2026} gives explicit actions of
all $\widehat{\mathrm{GT}}$ on the full profinite groups in types
$(g,0)$ and $(g,1)$; its introduction explicitly distinguishes these
actions from actions respecting the additional Lego rules.

The functorial actions in E.13 restrict to the universal-curve kernels
and preserve their peripheral twists.  Since genus one has the
congruence subgroup property, E.11 consequently yields an inclusion
$\widehat{\mathrm{GT}}\hookrightarrow\Gamma^{\mathrm{norm}}_{1,2}$
for that action.  The missing reverse inclusion in our reduction is the
assertion that every element of $A$ preserves the nonseparating inertia
classes in $E=\widehat P_{1,3}$.  The inertia-preserving automorphism
theorem does not itself prove this automatic-inertia assertion.

\statement{E.14}{Open}{The remaining automatic-inertia step in our elliptic reduction}
Every $a\in A$ from E.11 preserves the conjugacy classes of the closed
procyclic subgroups generated by nonseparating Dehn twists in
$\widehat P_{1,3}$.  If this holds, the inertia-preserving theorem
\citet[Theorem~A(ii)]{Boggi2026Automorphisms}, together with the
normalizer calculation for $\mathfrak S_2\subset\mathfrak S_3$, gives
$\Gamma^{\mathrm{norm}}_{1,2}\cong\widehat{\mathrm{GT}}$.
We leave this as the precise unproved step in this audit; we do not
claim that the general automatic-inertia problem has been settled.

For clarity, the comparisons obtained here are as follows.  Entries
describe the groups, not arbitrary choices of representatives of outer
actions.
\begin{center}
\begin{tabular}{L{.13\linewidth}L{.24\linewidth}L{.25\linewidth}L{.27\linewidth}}
Type & Revised arithmetic quotient & Early normalizer quotient & Related tower theorem\\ \hline
$(0,3)$ & $G_{\mathbb Q}$ & $\widehat{\mathrm{GT}}_0$ & pentagon extension required\\
$(0,4)$ & $G_{\mathbb Q}$ & $\widehat{\mathrm{GT}}$ & universal curve has type $(0,5)$\\
$(0,5)$ & $G_{\mathbb Q}$ & $\widehat{\mathrm{GT}}$ & stable genus-zero tower\\
$(1,1)$ & $G_{\mathbb Q}$ & $\widehat{\mathrm{GT}}$ & elliptic universal curve and E.10\\
$(1,2)$ & $G_{\mathbb Q}$ & $A/\mathfrak S_2$ as in E.11 & contains $\widehat{\mathrm{GT}}$ via E.13; E.14 unresolved here\\
\end{tabular}
\end{center}
The expected type-changing isomorphisms for the revised quotients are
all E.2.  For normalizer quotients, the genus-zero stable comparisons and
the comparison with $(1,1)$ above are group-theoretic theorems, whereas
equality with the arithmetic subgroup still requires E.8.  These are
three separate assertions.

\subsection{Arithmetic reconstruction uses more morphisms}

\statement{E.15}{Known}{What Pop's theorem reconstructs}
Let $U_0=\mathbb P^1\setminus\{0,1,\infty\}$ and let
$\mathcal V_{0,\mathrm{bir}}$ contain $U_0$, a cofinal system of open
subsets of $\mathcal M_{0,5}$, their inclusions and the restrictions of
the three functions $t_1,t_2,t_2-t_1$ to $U_0$ wherever defined.
Then the canonical map from $G_{\mathbb Q}$ to the compatible outer
automorphism group of the geometric fundamental-group functor of this
system is an isomorphism.  This is the application on pages~3--4 of
\citet{Pop2018Tripod}, from Theorem~2.7(1); its abelian-by-central
pro-$\ell$ counterpart is Theorem~2.6(1).
\begin{proof}
The category satisfies Hypothesis $(H)$ in Definition/Remark~2.2:
the variety has dimension greater than one, the designated functions
generate its function field over the algebraic closure, and the specified
birational system is rigid.  Theorem~2.7 identifies the automorphisms
of the tame fundamental-group functor; in characteristic zero this is
the full geometric \emph{profinite} functor.  The category has infinitely
many localization objects even though only three rational functions are
designated.  Removing those objects destroys the stated hypothesis.
Thus this theorem cannot be specialized to just
$(\mathcal M_{0,4},\mathcal M_{0,5})$ while keeping its conclusion.
\end{proof}

\statement{E.16}{Known}{The 2025 combinatorial reconstruction is not arithmetic surjectivity onto GT}
\citet[Theorem~C(ii)]{HoshiMochizukiTsujimura2025} reconstruct the
$\widehat{\mathrm{GT}}$-conjugacy class of $G_{\mathbb Q}$ as the maximal
closed subgroups satisfying their AA conditions.  The construction uses
the abstract configuration group
$\pi_1(\operatorname{Conf}_2(U_{0,\overline{\mathbb Q}}))$.
The additional COF, RGC, arithmetic Belyi compatibility and ampleness
requirements are indispensable hypotheses, defined in their
Definitions~3.3 and~5.12.  Their Theorem~A constructs field operations;
Theorem~B proves an injectivity assertion under additional hypotheses.
None states that all of $\widehat{\mathrm{GT}}$ has these properties.

We located this result through \citet[\S4]{Collas2026}, and then opened
the original, dated 14 March 2025.  A characterization of a conjugacy
class of arithmetic subgroups answers one formulation of scan~479's
request for an algebraic description.  It is not the equality in E.8,
and does not classify all admissible splittings of the moduli extensions.

\statement{E.17}{Proved}{The modular consequence of the claimed section theorem}
Assume the section theorem for $X_1(N)$ claimed in Corollary~8.5 of
the 2026 manuscript in family~019, with $g(X_1(N))\geq2$.
Let a section of the arithmetic extension for type $(1,1)$ lift to the
open subgroup belonging to the fine modular curve $Y_1(N)$.
Its composite with
$\pi_1(Y_1(N))\to\pi_1(X_1(N))$ is, up to geometric conjugacy, the
section of a rational cusp of $X_1(N)$.
\begin{proof}
The composite is a section because both maps are over $G_{\mathbb Q}$.
The asserted proper-curve section theorem makes it the section of a
rational point.  Mazur's theorem implies that all rational points of
$X_1(N)$ are cusps when its genus is positive; see
\citet[Introduction, Theorems~(7),(7')]{Mazur1977}, and the detailed
conditional argument in Part D, Statement D.10.  The level assumption
here also removes the generic elliptic stabilizer, so that the subgroup
is the fundamental group of a fine curve, not its coarse stack quotient.
\end{proof}
This does not show that the original affine section is conjugate into
a cusp decomposition group.  Nor does it show nonexistence: tangential
sections at rational cusps supply affine sections.  The unrefereed
proper-curve theorem does not treat affine curves, moduli stacks or all
splittings of $N^{\mathrm{arith}}_{1,1}\to G_{\mathbb Q}$.
The hypothesis and the status of family~019 are retained in the results
register; no unconditional section theorem is asserted here.

\statement{E.18}{Known}{Arithmetic rigidity is a different automorphism problem}
For a number field $k$, an odd prime $\ell$, and $m\geq4$,
\citet[Theorem~A]{Nakamura1994} identifies the $G_k$-compatible
automorphisms of the geometrically pro-$\ell$ arithmetic fundamental
group of $\mathcal M_{0,m}$, modulo geometric inner automorphisms, with
$\operatorname{Aut}_k(\mathcal M_{0,m})$.  The centralizer of the
geometric outer Galois image is $\mathfrak S_3$ at $m=4$, and
$\mathfrak S_m$ at $m\geq5$.  This has an \emph{odd-prime} hypothesis;
it is not a theorem that the unrestricted outer automorphism group is
finite.  The arithmetic compatibility condition cannot be deleted to
deduce E.8.  His Theorem~B is the corresponding Lie-algebra centralizer
statement for $m\geq5$.  This is the precise rigidity result used here,
rather than an unspecified assertion of ``Galois rigidity''.

\subsection{The rational comparison and the maps to the pro-\texorpdfstring{$\ell$}{l} tower}

For a field $k$ of characteristic zero, unipotent completion means the
universal pro-unipotent affine $k$-group receiving the discrete group,
or its continuous $\ell$-adic variant when $k=\mathbb Q_\ell$.
It is not the discrete fundamental group retained by arbitrary locally
constant $D(\mathbb Q)$-coefficients, and it is not the entire
pro-algebraic completion.  In the braid tower we complete the pure braid
groups and their path torsors, retaining the permutations and the
cabling, deletion and reassociation operations.  Completing the whole
mapping class group indiscriminately would define a different object.
The grading below is positive word weight; motivic cohomological weight
is its negative double.

Let $\mathfrak t_j$ be the infinitesimal braid Lie algebra on the
$t_{ab}=t_{ba}$, with
$[t_{ab},t_{cd}]=0$ for distinct indices and
$[t_{ab},t_{ac}+t_{bc}]=0$.  Let
$W=\bigoplus_{w\geq1}W_w\subseteq\operatorname{Lie}_{\mathbb Q}
\langle x,y\rangle$ be the solution space
\begin{align*}
 \psi(x,y)+\psi(y,x)&=0,\\
 \psi(x,y)+\psi(y,-x-y)+\psi(-x-y,x)&=0,\\
 \psi(t_{12},t_{23}+t_{24})+\psi(t_{13}+t_{23},t_{34})
 &=\psi(t_{23},t_{34})+\psi(t_{12}+t_{13},t_{24}+t_{34})
    +\psi(t_{12},t_{23}).
\end{align*}
The third equation lies in $\mathfrak t_4$.  The bracket is
\[
 \{\psi,\varphi\}=D_\psi\varphi-D_\varphi\psi+[\psi,\varphi],
 \qquad D_\psi(x)=0,\quad D_\psi(y)=[y,\psi].
\]
The ordinary commutator on the free Lie algebra is not this bracket.

\statement{E.19}{Known}{The precise rational source, target and established injection}
The weight-completed Lie algebra of $\mathrm{GRT}_1$ is
$\widehat W=\prod_w W_w$ with the Ihara bracket, and
$\mathrm{GRT}=\mathbb G_m\ltimes\mathrm{GRT}_1$.
It is the graded automorphism group for the completed parenthesized
chord tower; the corresponding filtered braid tower has automorphism
group $\mathrm{GT}(k)$.  A choice of a $k$-associator identifies these
two groups, noncanonically.  The symmetry and hexagon occur at the
four-point level and the pentagon at the five-point level; the same
operations determine the action at every genus-zero level.

For the de Rham fibre functor on $\mathrm{MT}(\mathbb Z)$, let
$G_{\mathrm{MT}}=\mathbb G_m\ltimes U_{\mathrm{MT}}$ and let
$\mathfrak u_{\mathrm{MT}}^{\mathrm{gr}}$ denote its weight-graded
unipotent Lie algebra, reindexed by positive word weight.  There is an
injective graded Lie homomorphism
\[
 j:\mathfrak u_{\mathrm{MT}}^{\mathrm{gr}}
   \cong\operatorname{Lie}_{\mathbb Q}\langle e_3,e_5,e_7,\ldots\rangle
   \hookrightarrow W.
\]
It integrates after weight completion to a monomorphism of
pro-unipotent groups.  The generators themselves are not canonical.
\begin{proof}
Drinfeld's definitions in \S\S5--6 and Propositions~5.1 and~5.5
give the associator torsor and the two actions
\citep{Drinfeld1990}.  The structure-preserving categorical formulation
is \citet[Proposition~3.4, Theorems~1--3, Corollary~4.1]{BarNatan1998};
its operations must be retained when invoking that formulation.
Differentiating gives exactly the three equations displayed above.
\citet[Theorem~1.1 and its stated consequence, Conjecture~1]{Brown2012}
proves faithfulness of the mixed Tate action.  The containment in the
Grothendieck--Teichm\"uller solution Lie algebra is stated explicitly in
\citet[\S1.4.2]{Brown2021Depth}, using Brown's theorem and Drinfeld's
equations.  The free odd-weight source is the structure of
$\mathrm{MT}(\mathbb Z)$ recalled in \citet[\S1.1]{Brown2012}.
In characteristic zero, exponential and logarithm give the equivalence
between finite-dimensional nilpotent Lie algebras and unipotent groups;
passing through finite weight quotients gives the completed assertion.
\end{proof}

\statement{E.20}{Proved}{The exact conditional rational counterpart}
Assume the all-weight bound of family~008, Corollary~5.4:
\[
 \dim_{\mathbb Q} W_w\leq
 \dim_{\mathbb Q}\operatorname{Lie}\langle e_3,e_5,\ldots\rangle_w
 \quad\hbox{for every }w.
\]
Then the map $j$ of E.19 is an isomorphism in every weight and after
weight completion.  It gives
$U_{\mathrm{MT}}\cong\mathrm{GRT}_1$ and, with the matching Tate
character and grading, $G_{\mathrm{MT}}\cong\mathrm{GRT}$.
Via a chosen associator this is the analogous statement for
$\mathrm{GT}$ on the unipotent genus-zero braid tower.
\begin{proof}
The injection $j_w$ has a source of the dimension on the right.
The assumed opposite inequality forces equality and surjectivity.
Every weight piece is finite dimensional, so the product of these
isomorphisms is an isomorphism for the complete weight filtration;
the bracket in each weight involves finitely many lower weights.
Exponentiation on the finite nilpotent quotients and inverse limit
give the group isomorphism.  The scalar actions agree by their weight
definitions (with the inverse character if the negative motivic-degree
convention is used), so they extend over $\mathbb G_m$.
\end{proof}
The conclusion is \emph{claimed (OpenAI, 2026, unrefereed)} when
Corollary~5.4 is supplied by that manuscript.  The implication just
proved is unconditional logic; its all-weight premise is not certified
by our bounded calculation or by reading a comparator statement.

\statement{E.21}{Known}{The realization chain and what its arrows do}
Fix a prime $\ell$ and compatible tangential markings.  Write
$\mathrm{GT}_{\ell}$ for the group of invertible pairs satisfying the
pro-$\ell$ versions of Drinfeld's equations, with
$\lambda\in\mathbb Z_\ell^{\times}$.  This is generally not a
pro-$\ell$ group, because its cyclotomic quotient need not be one.
There are natural maps
\[
 G_{\mathbb Q}\hookrightarrow\widehat{\mathrm{GT}}
   \longrightarrow\mathrm{GT}_{\ell}
   \hookrightarrow\mathrm{GT}(\mathbb Q_\ell).
\]
The first arrow is a continuous closed embedding.  The second is
continuous and is not injective; we assert no surjectivity.
The third is injective and continuous for the topology on finite weight
quotients.  Neither its image nor the image of $G_{\mathbb Q}$ is
surjective, topologically dense or open in
$\mathrm{GT}(\mathbb Q_\ell)$ with this inverse-limit topology.
Surjectivity or openness of the arithmetic image \emph{inside}
$\mathrm{GT}_{\ell}$ is not concluded here.
\begin{proof}
Drinfeld constructs the completion comparison maps in \S4
\citep{Drinfeld1990}; the continuous pro-$\ell$ to unipotent
comparison is also \citet[Theorem~A.6 and Corollary~A.10]{HainMatsumoto2003}.
Faithful Magnus expansion of the free pro-$\ell$ group gives
injectivity of its map to the characteristic-zero completed free
unipotent group, hence injectivity of the pair map.
The arithmetic pro-$\ell$ action factors through the Galois group
$G_{\mathbb Q,\{\ell\}}$ of extensions unramified outside $\ell$;
this is the factorization recalled and used in
\citet[Introduction and \S8.4]{HainMatsumoto2003}.
Nontrivial inertia at a prime different from $\ell$ therefore lies in
the kernel of the composite.  Belyi injectivity in the first arrow
forces a nontrivial kernel in the second.

The pro-$\ell$ solution group is compact.  Its scalar image lies in
$\mathbb Z_\ell^{\times}$, a proper closed subset of
$\mathbb Q_\ell^{\times}$, proving nonsurjectivity and failure of
density in the full rational group.  Its image is compact and hence
closed.  It cannot be open: any neighbourhood of the identity in the
inverse-limit weight topology contains the kernel of a sufficiently
large finite weight projection.  By Brown's independent odd-weight
elements, such a kernel contains a closed one-parameter subgroup
isomorphic to the additive group $\mathbb Q_\ell$ in an arbitrarily
large weight, and is not compact.  A compact subgroup cannot contain
that closed noncompact subgroup.  The same argument applies to the
compact arithmetic image.
\end{proof}

Hain--Matsumoto is a statement about weighted completion and the
\emph{image} of an arithmetic action, not faithfulness of
$G_{\mathbb Q}\to\mathrm{GT}_{\ell}$.  More precisely their
Theorem~7.3 computes the free graded Lie algebra of the weighted
completion of $G_{\mathbb Q,\{\ell\}}$ with respect to the cyclotomic
character; this initially includes weight one.  Their Theorem~8.4
identifies the graded Lie image with the rationalized graded image of
the Galois action.  Lemma~8.5 eliminates weight one in the outer action;
Remark~8.6 identifies the nonzero odd abelian classes by Soul\'e's
characters.  Their Theorem~9.8 gives the constrained mixed Tate
factorization.  Brown's faithfulness then supplies the free motivic
image, while E.20, conditionally on the upper bound, says that this
image exhausts $W\otimes\mathbb Q_\ell$.  The resulting density is
\emph{Zariski density in each finite-dimensional algebraic quotient}.
It is compatible with the preceding failure of topological density in
the full group of $\mathbb Q_\ell$-points.  No argument here divides by
$2$ or $3$ in a profinite inertia group; the all-prime Hain--Matsumoto
construction and the separate characteristic-two lattice calculation
must not be confused with Nakamura's odd-prime theorem E.18.

\statement{E.22}{Proved}{The needed profinite analogue of the upper bound}
For the closed embedding $\iota:G_{\mathbb Q}\hookrightarrow
\widehat{\mathrm{GT}}$, the following are equivalent:
\begin{enumerate}
\item $\iota$ is surjective;
\item its image is topologically dense;
\item for every open normal subgroup $V$ of
$\widehat{\mathrm{GT}}$, the map
$G_{\mathbb Q}\to\widehat{\mathrm{GT}}/V$ is surjective.
\end{enumerate}
This finite-quotient exhaustion, not a rational dimension inequality,
is the required profinite counterpart of Corollary~5.4.
\begin{proof}
A compact image in a Hausdorff profinite group is closed.  Density is
therefore equivalent to equality.  Density implies surjectivity on
every finite discrete quotient.  Conversely, if a closed subgroup $H$
is proper, a point outside $H$ has an open neighbourhood disjoint from
$H$; shrinking to a coset of an open normal subgroup $V$ gives a coset
missed by $H$ in the finite quotient.  This contradicts the third
condition.
\end{proof}

\statement{E.23}{Refuted}{Equality of unipotent completions does not imply profinite equality}
The inference ``a subgroup induces an isomorphism of characteristic-zero
unipotent completions, therefore it is the whole profinite group'' is
false, already for $1\subset C_2$.
\begin{proof}
Every finite-order element of a characteristic-zero unipotent group is
the identity.  Indeed, write it as $\exp X$ in a faithful upper
unitriangular representation; $(\exp X)^r=1$ implies $rX=0$, hence
$X=0$.  Therefore every homomorphism $C_2\to U$ is trivial for every
unipotent $U$, and the unipotent completions of $1$ and $C_2$ are both
trivial.  Their profinite completions are themselves, and differ.
The supplied program checks the two-element multiplication table and
its associativity; the torsion-free argument is the proof of the
completion assertion.  This counterexample does not say that the
arithmetic inclusion is proper.  It refutes only the proposed transfer
principle.
\end{proof}

\statement{E.24}{Proved}{Conditional abelianizations of GRT}
Assume the freeness assertion of family~008, Theorem~1.1.  Then
\[
 W^{\mathrm{ab}}\cong\bigoplus_{k\geq1}\mathbb Q e_{2k+1},
 \qquad
 (\mathrm{GRT}_1)^{\mathrm{ab}}
       \cong\prod_{k\geq1}\mathbb G_a(2k+1),
 \qquad
 (\mathrm{GRT})^{\mathrm{ab}}\cong\mathbb G_m.
\]
The middle and last abelianizations are in the category of affine
pro-algebraic groups, equivalently by closed commutator subgroups.
The last formula, with the scalar character, is the rational analogue
of Grothendieck's cyclotomic abelianization statement; the middle
formula is not.
\begin{proof}
The abelianization of a free graded Lie algebra is its generating
space.  Passing to the weight completion replaces the direct sum by
the product, and exponentiation identifies an abelian unipotent group
with its additive Lie algebra.  In the full semidirect product, every
positive-weight abelian piece has commutator action
$(h-1)v=(2^{\mathrm{wt}(v)}-1)v$ for the scalar $h=2$.
The coefficient is nonzero and invertible over $\mathbb Q$.
Thus all unipotent weight pieces vanish in the abelianization.
This argument can be applied successively in every finite weight
quotient; inverse limit gives the closed-commutator assertion.  The
remaining quotient is the already abelian group $\mathbb G_m$.
No assertion about the abstract, nonclosed derived subgroup of an
infinite product is being made.
\end{proof}

\statement{E.25}{Known}{The arithmetic meaning of the odd generators}
On the motivic side the odd generating lines are dual to
\[
 \operatorname{Ext}^1_{\mathrm{MT}(\mathbb Z)}
       (\mathbb Q(0),\mathbb Q(2k+1))
       \cong K_{4k+1}(\mathbb Z)\otimes\mathbb Q,
       \qquad k\geq1.
\]
Each line has dimension one.  Their $\ell$-adic realizations are
measured by the nonzero Soul\'e characters in weight $2k+1$.
In particular they are abelian classes in the unipotent or
cyclotomic-kernel Lie algebra, not extra ordinary abelian characters
of $G_{\mathbb Q}$ contrary to Kronecker--Weber.
\begin{proof}
This is the mixed Tate description recalled in
\citet[\S1.1]{Brown2012}, together with
\citet[Theorem~7.3 and Remark~8.6]{HainMatsumoto2003}.
The $G_{\mathbb Q}$-equivariant abelianization in the latter is
$\bigoplus_{k\geq1}\mathbb Q_\ell(2k+1)$.
Taking the full semidirect product with the cyclotomic torus kills
these positive-weight coinvariants, exactly as in E.24.
\end{proof}

\subsection{Audit of the new all-weight upper bound}

\statement{E.26}{Open}{The imported 2026 claim and the extent of our verification}
Family~008, \citet[Theorem~1.1 and Corollary~5.4]{OpenAI2026DD},
claims that $W$ is free on one generator in every odd weight at least
three.  We read \S\S2--5, including the all-weight reduction, and
verified the implication E.20 using the established Brown lower bound.
The following audit records the actual checks and their limits.
We have not independently certified every all-weight rewriting and
deletion argument, so our mathematical use of the upper bound remains
conditional.  ``Open'' here records that outstanding verification in
this deliverable; it does not erase the announced proofs described
below.

\paragraph{Integral reduction (\S2).}
The use of $R=\mathbb Z_{(2)}$ is appropriate.  In a fixed weight the
free ordinary Lie lattice is finite free by its integral Hall--Lyndon
basis.  The solution lattice
$L_w=W_w\cap\operatorname{Lie}_R\langle x,y\rangle_w$ is saturated:
if $2a\in L_w$, then $a\in W_w$ because the equations are rational
linear, and $a$ already lies in the ambient lattice.  Hence its reduction
embeds in the ambient reduction, and
$\dim_{\mathbb F_2}\overline L_w=\dim_{\mathbb Q}W_w$.
We followed this proof of Proposition~2.5.  It does \emph{not} assert
that $\overline L_w$ is the entire mod-two solution space.
The degree-one pentagon leaves $-\alpha t_{12}-\beta t_{34}$,
and the degree-two three-cycle gives $3[x,y]$; thus Lemma~2.1 indeed
gives $W_1=W_2=0$.

Lemma~2.3 uses the standard free-fibre decomposition of the
infinitesimal braid Lie algebra; its explicit derivations give the
required integral torsion-freeness of the targets.  We followed the
presentation comparison.  For Lemma~2.4 we followed the reduction of
the pentagon in the free fibre, the centralizer argument, and the
linearization giving the special-derivation identity and its transverse
derivative.  Its division by the nonzero integer $w+1$ takes place
over $\mathbb Q$ \emph{before} reducing the saturated lattice; it is
not a division by an even integer in $\mathbb F_2$.  This distinction
is necessary for the stated proof to make sense.

\paragraph{Filtered comparison (\S3).}
The substitution
$A=x^2$, $C=[x,y]$, $B=y$ is made in the associative algebra in
characteristic two, with the number of occurrences of $B$ as filtration.
The prefix words $xx,xy,y$ give the asserted injective associative
encoding.  The source and target commutator tables have a common
count-preserving part: correction terms raise the count in the braid
source and lower it in the cyclotomic arrangement target.  The ordering
rule replaces a base letter followed by a fibre letter by the opposite
order plus its commutator, decreasing the stated inversion measure.
The proof needs independent normal words in the \emph{source}; mere
spanning in the target is enough.  We checked this direction of the
argument and selected entries of the tables (3.6)--(3.12).  We did not
rederive every entry or mechanize the complete induction.  Proposition~3.2
is therefore still an imported all-weight step, not a consequence of
our finite computation.

\paragraph{Deletion representation and leading-term injectivity (\S4).}
The nine type-$B_3$ hyperplanes are first defined over $\mathbb Q$.
The plus and minus hyperplanes remain distinct symbols after mod-two
reduction; merging them geometrically in characteristic two would change
the arrangement Lie algebra.  The deletion operator at a letter has
coefficient equal to the valuation on its right edge minus the
valuation on its left edge, with equal-letter differences assigned zero.
We followed Proposition~4.1's logarithmic flatness and residue argument
over the integers.  Our independent implementation of those edge
valuations tests all 13 codimension-two flats, every incident distinguished
hyperplane, and every word of length at most five in the three-letter
alphabet.  All 13,104 integer commutator instances vanish.
This checks signs, hyperplane labels and the repeated-letter convention
in that range.  It does not prove the representation in unbounded length.

We followed how Lemma~4.3 uses the minimal number of $A$'s, terminal
deletion and reversal to constrain a leading polynomial after $A=0$.
Proposition~4.4 then uses the reduced special identity to compare
coefficients, separating the cases of at least two $A/C$ letters and
the small boundary cases.  In particular the $r=2$ expression
$ABB+BBA$ is not the uniform coefficient sum; the proof cannot discard
that case by an informal characteristic-two symmetry argument.
We checked the displayed reductions and case division, but did not
independently reconstruct the general coefficient-exchange induction.

\paragraph{The image bound (\S5).}
The even-index bound is an application of free-Lie elimination.
Lemma~5.2 uses the Ore-extension calculation with
$D(B)=C$, $D(C)=0$, $D^2=0$, giving invariance under
$B\mapsto B+s(C)$.  In Proposition~5.3 the Vandermonde variables
are independent indeterminates; the argument takes place over a
polynomial ring, not by evaluating a supposedly invertible matrix on
the two elements of $\mathbb F_2$.  Divisibility by the successive
differences in the encoding excludes the index-zero elimination
generator.  Lazard elimination then leaves the ordinary free Lie
algebra generated by
$g_k=\operatorname{ad}_C^k(B)$ for $k\geq1$, whose weights are
$2k+1$.  We followed the polynomial divisibility and ordinary-Lie
retraction steps.  The distinction between ordinary Lie elements and
all primitive elements in characteristic two is consistently necessary.
Combining this image bound with the injective projection and the finite
filtration yields Corollary~5.4 exactly as stated.  Proposition~5.5's
leading-bracket calculation is not needed for E.20, since Brown already
provides the compatible Lie injection.

\paragraph{What this audit establishes.}
The integral dimension transfer, low-weight vanishing, conditional
dimension squeeze, completion passage, abelianization calculation and
bounded deletion check have been independently checked here.  The
unbounded commutator-table/normal-form comparison of \S3 and the
coefficient-detection induction of \S4 remain the principal steps not
fully rederived by us.  We did not use the alternative lower-bound
construction in \S\S6--8, its regularized-transport proposition, or
the graph-complex corollary.  They are claimed, steps not checked by us.
The program additionally computes, from the free-Lie generating series,
the predicted weight dimensions through 30.  Those numbers describe
the proposed bound; they are not a computation of $\dim W_w$.

\statement{E.27}{Refuted}{The attached comparator is not a checked Lean proof}
The claim that the attached file
\texttt{lean/DeligneDrinfeld.lean} itself provides a kernel-checked
proof of Theorem~1.1 is false.
\begin{proof}
We read the supplied file.  It defines the solution equations, Ihara
bracket and the claimed graded and completed statement, and ends with
\texttt{theorem main : MainStatement := by sorry}.
The scope note describes a formalization elsewhere, but that description
does not remove the placeholder in this comparator.  We did not fetch,
build or audit a separate Lean proof repository.  Thus this attachment
confirms the shape of the target statement only.  This is a refutation
of a certification claim, not of the mathematical theorem.
\end{proof}

\statement{E.28}{Known}{The status search finds subsequent proof claims}
As of our search on 11 November 2026,
\citet[Theorem~1.1]{Willwacher2026After} presents a simplified proof
of the Deligne--Drinfeld conclusion following family~008; the opened
version is arXiv:2610.15592v4, dated 28 October 2026.
We read its principal statements and reduction, not a complete
independent proof audit.  It is an additional primary preprint source,
so it would be inaccurate to report that no subsequent mathematical
check or proof treatment exists.  Its existence does not turn the
attached comparator into a checked formal proof, and it does not
address the profinite equality E.8.  Our conditional labels for use of
the supplied upper bound remain explicit.

\subsection{Consequences for the requested categorical formulation}

The group acting on a completed degeneration category in Part F must
be specified by its retained operations.  For the marked genus-zero
system with common inertia exponent, relabellings, hole-filling and
clutching, its restriction is the group in E.4.  A proof that these
restrictions determine and lift category automorphisms is additional
work; it is not part of the wording of the Harbater--Schneps theorem.
Part F supplies that comparison.  Once it is supplied, replacing each
pure braid group and tangential path torsor by its unipotent completion
and retaining cabling and reassociation gives the filtered
$\mathrm{GT}$-tower, or, after an associator, the graded
$\mathrm{GRT}$-tower in E.19.  Its motivic arithmetic equality is
precisely the conditional theorem E.20.  The revised profinite
arithmetic quotient is already $G_{\mathbb Q}$ by E.2; equality with
the full geometric tower automorphism group still asks E.8.

Neither the two excluded low-dimensional elliptic and rational
nonhyperbolic types nor orientation-reversing mapping classes can be
inserted into these assertions without changing the hypotheses.
In particular, quotienting by orientation reversal would kill complex
conjugation in a real model.  The orientation-positive geometric
subgroup in E.1--E.2 is therefore essential.  Stack stabilizers at
elliptic curves with extra automorphisms, including those of orders
$4$ and $6$, are retained throughout; passing to a coarse $j$-line
would not preserve the fundamental-group extensions used above.


